\subsection{The exact weight-seven image}
\label{subsec:weight-seven-exact}

Two coefficients of the full pentagon, together with the signed
congruence identity, determine the exact image. The coefficient calculation uses the terms of
weights eight and nine arising from the ordinary pure-braid action.

\begin{lemma}[A matrix lower-central estimate]
\label{lem:weight-seven-matrix-filtration}
Let $Q$ be the closed subgroup of $\mathrm{SL}_2(\mathbf Z_2)$
generated by $U^2,V^2$, where
\[
 U=\begin{pmatrix}1&1\\0&1\end{pmatrix},\qquad
 V=\begin{pmatrix}1&0\\-1&1\end{pmatrix}.
\]
Let $\mathcal B_r$ be the principal congruence group modulo $2^r$,
and let $\mathcal A_r$ consist of determinant-one matrices whose
diagonal entries are $1$ modulo $2^r$ and whose off-diagonal
entries are zero modulo $2^{r+1}$. Then
\[
 \gamma_{2k}(Q)\subseteq\mathcal A_{3k-1},\qquad
 \gamma_{2k+1}(Q)\subseteq\mathcal B_{3k+1}\quad(k\ge1).
\]
In particular, $\gamma_7(Q)\subseteq\mathcal B_{10}$,
$\gamma_9(Q)\subseteq\mathcal B_{13}$ and
$\gamma_{10}(Q)\subseteq\mathcal A_{14}$.
\end{lemma}

\begin{proof}
For $M=(\alpha,\beta;\gamma,\delta)$ of determinant one,
\[
 [M,U^2]=
 \begin{pmatrix}
 1-2\alpha\gamma&2(\alpha^2-1)+4\alpha\gamma\\
 -2\gamma^2&1+2\alpha\gamma+4\gamma^2
 \end{pmatrix}.
\]
It follows for $r\ge1$ that
$[\mathcal B_r,U^2]\subseteq\mathcal A_{r+1}$ and
$[\mathcal A_r,U^2]\subseteq\mathcal B_{r+2}$.
Exchanging the coordinates gives the same assertions for $V^2$.
The determinant identity also gives
$\alpha-\delta\equiv0\pmod{2^{r+1}}$ in $\mathcal A_r$.
The conjugation formulas therefore show that $U^2,V^2$ normalize
$\mathcal A_r$; they also normalize $\mathcal B_r$.
Hence the two generator calculations imply
\[
 [\mathcal B_r,Q]\subseteq\mathcal A_{r+1},\qquad
 [\mathcal A_r,Q]\subseteq\mathcal B_{r+2}.
\]
Starting with $Q\subseteq\mathcal B_1$ and iterating proves the
stated lower-central estimates.
\end{proof}

\begin{theorem}\label{thm:weight-seven-exact-image}
One has
\[
 m_7(D_5)=m_7(J_5)=128\mathbf Z_2.
\]
Consequently, for $D_7=D_5\cap\ker m_7$ and $J_7=D_7\cap H$,
\[
 D_5=D_7J_5,\qquad G=D_7H,\qquad
 B\simeq D_7/\overline{[G,D_7]J_7}.
\]
The subgroup $D_7$ is proper, closed and normal in $G$, and
\[
 \eta_7:D_5\longrightarrow\mathbf F_2,
 \qquad \eta_7(g)=m_7(g)/128\pmod2
\]
is a continuous $G$-invariant character with nonzero arithmetic
restriction.
\end{theorem}

\begin{proof}
It remains to prove $m_7(D_5)\subseteq128\mathbf Z_2$, since an
arithmetic value of exact valuation seven was already constructed.
Use the integral degree-seven Hall coordinates $a,b$, for which
the leading pentagon gives $a=16A$, $b=19A$ with $A\in\mathbf Z_2$.
Define
\[
 M_f=2^{-10}\bigl(f(U^2,V^2)\bigr)_{12}\pmod{16}.
\]
The matrix lower-central estimates give
$\gamma_7(Q)\subseteq\mathcal B_{10}$ and
$\gamma_{10}(Q)\subseteq\mathcal A_{14}$.
Thus $M_f$ is an integral linear form modulo sixteen in the Hall
coordinates of weights seven through nine. Higher factors are
invisible modulo $2^{14}$, and products and powers of the remaining
factors are linear at that precision.

Let $P_\alpha$ and $P_\beta$ be the full-pentagon coefficient forms
defined below. They are integral and vanish for every GT
word. The finite coefficient identity is
\begin{equation}\label{eq:weight-seven-upper128-certificate}
 6(a+b)-8P_\alpha-4P_\beta-5M_f
       \equiv\sum_r e_rE_r\pmod{16},
\end{equation}
where the $E_r$ are coefficients of the two-cycle and three-cycle
relations through degree eleven. All these forms, and every
coefficient $e_r$, are specified below.

For an element of $D_5$, the defining equations make
$P_\alpha=P_\beta=E_r=0$. The normalized signed congruence identity
makes $M_f=0$. Hence $6(a+b)\equiv0\pmod{16}$, so
$a+b\equiv0\pmod8$. Since $a+b=35A$ and $35$ is odd, $A$ is
divisible by eight. Therefore $a\in128\mathbf Z_2$.

The arithmetic construction and compactness now give equality of
both images. The scalar conjugation law proves normality and
invariance, while the arithmetic value proves properness and
nonvanishing. Equality of the two images gives $D_5=D_7J_5$ and
then $G=D_7H$. The relative quotient identification follows from
the preceding centralization argument.
\end{proof}

\paragraph{The full-pentagon coefficient forms.}
Let $\Sigma_i$ be the ordinary Artin automorphism of the free group
on $x_1,x_2,x_3,x_4$:
\[
 \Sigma_i(x_i)=x_ix_{i+1}x_i^{-1},\qquad
 \Sigma_i(x_{i+1})=x_i,
\]
with the other generators fixed. Its inverse sends $x_i$ to
$x_{i+1}$ and $x_{i+1}$ to $x_{i+1}^{-1}x_ix_{i+1}$.
Put
\[
 \mathcal T_{ij}=(\Sigma_{j-1}\cdots\Sigma_{i+1})
                 \Sigma_i^2
                 (\Sigma_{j-1}\cdots\Sigma_{i+1})^{-1}.
\]
These are the images of the pure generators in the pentagon
convention of \cite[Section~5, equation~(17) and the following displayed
formula]{Enriquez2007}; see also \cite[Section~1]{NS2000}.
Use the following five pairs and signs:
\[
\begin{array}{c|c|c|r}
i&\mathcal U_i&\mathcal V_i&\varepsilon_i\\ \hline
1&\mathcal T_{12}&\mathcal T_{23}\mathcal T_{24}&1\\
2&\mathcal T_{13}\mathcal T_{23}&\mathcal T_{34}&1\\
3&\mathcal T_{23}&\mathcal T_{34}&-1\\
4&\mathcal T_{12}\mathcal T_{13}&\mathcal T_{24}\mathcal T_{34}&-1\\
5&\mathcal T_{12}&\mathcal T_{23}&-1.
\end{array}
\]
Products mean composition of automorphisms in the written order.
Write $x_j=1+X_j$ and let $\Delta_{\mathcal U}=\mathcal U-\mathrm{Id}$
act linearly on the completed noncommutative Magnus algebra.
Every pure braid acts trivially in degree one, so each $\Delta$
raises degree by at least one.

For a standard Lyndon group word, let
$Q_w(S,T)=L_w(1+S,1+T)-1$, truncated after degree nine. Define
\begin{equation}\label{eq:weight-seven-pentagon-functional}
 P_\xi(h)=[\xi]\sum_{i=1}^5\varepsilon_i
       \sum_{7\le |w|\le9}h_w
       Q_w(\Delta_{\mathcal U_i},\Delta_{\mathcal V_i})(X_1),
\end{equation}
for the two words
\[
 \alpha=X_1^6X_2X_3X_2X_4,
 \qquad
 \beta=X_1^4X_2X_3X_2X_4X_2X_4.
\]
Multiplication in the polynomial $Q_w$ means composition of the
operators. Formula~\eqref{eq:weight-seven-pentagon-functional}
defines integral linear forms in the Hall coordinates.

We verify their precise scope. A Hall factor of weight at least
seven has Artin operator equal to the identity through degree
increase six. Products of two nonconstant such factors have
degree increase at least fourteen. Thus, through output degree
ten on $X_1$, the group pentagon is exactly the signed sum in
\eqref{eq:weight-seven-pentagon-functional}, linear in the Hall
coordinates. Hall weights at least ten have no effect at this
output degree. Moreover, every operator monomial that can contribute
has at least seven factors $\Delta$. A piece of one factor that
raises degree by four would give total degree increase at least
ten, hence output degree at least eleven. It therefore suffices
to expand the four base Artin images through degree four and use
their full substitution operators at the required precision.
This retains all terms needed from the full pure-braid action.
Since the full pentagon is an identity, $P_\alpha=P_\beta=0$
for the word of every element under consideration.

\paragraph{Complete finite data for the congruence.}
Form the integer Magnus polynomials through degree eleven, using
$(1+A)^{-1}=\sum_{j=0}^{11}(-A)^j$ and the standard group
commutator recursion. The $389$ Hall columns have weights seven
through eleven. Set
\[
 F=1+\sum_{7\le |w|\le11}h_w(L_w-1),\qquad Z=Y^{-1}X^{-1},
\]
and define
\[
 I_v=[v]\bigl(F(Y,X)+F(X,Y)-2\bigr),\qquad
 II_v=[v]\bigl(F(Z,X)+F(Y,Z)+F(X,Y)-3\bigr).
\]
They are linear, because products of two nonconstant factors begin
in degree fourteen. To compute $M_f$, form the two-by-two matrix
words $P_x=U^2$, $P_y=V^2$,
$P_w=P_uP_vP_u^{-1}P_v^{-1}$ modulo $2^{14}$; its coefficient on
$h_w$ is $(P_w)_{12}/2^{10}$ modulo sixteen. The higher Hall
columns have zero contribution as above. The two pentagon forms
are given by \eqref{eq:weight-seven-pentagon-functional}, with
the same Hall column order.

For the following table, encode $I_v$ by
$r=2^{|v|}+\operatorname{bin}(v)$ and $II_v$ by
$r=4096+2^{|v|}+\operatorname{bin}(v)$, with $S=0$, $T=1$.
The complete $244$ nonzero coefficients $e_r$ in
\eqref{eq:weight-seven-upper128-certificate} are listed below.
The asserted congruence is an identity of integral linear forms
on all $389$ Hall columns, so every entry is checkable from the
displayed finite recursions.


{\small
\[
\begin{array}{rr@{\qquad}rr@{\qquad}rr@{\qquad}rr}
r&e_r&r&e_r&r&e_r&r&e_r\\ \hline
129&-2&258&-1&259&3&261&4\\
262&-2&268&2&518&-6&519&-1\\
521&-6&523&6&524&-5&526&-3\\
534&7&1036&-1&1038&-5&1039&-4\\
1042&-5&1044&1&1047&-4&1051&-4\\
1052&-7&1054&-4&1057&-5&1058&-3\\
1059&-5&1063&8&1067&8&1068&5\\
1069&8&1070&-5&1074&-6&1076&8\\
1084&-4&1086&4&1094&-3&1107&4\\
1108&-6&1109&-4&1118&-4&1123&4\\
1144&-7&1180&-6&1208&-2&1212&4\\
1280&-4&1584&-4&2072&-3&2076&-7\\
2078&-4&2079&-4&2084&-6&2086&4\\
2087&8&2089&-5&2090&-5&2091&4\\
2094&4&2096&5&2097&-3&2098&-1\\
2100&-6&2101&8&2103&-4&2108&-4\\
2110&-4&2114&-1&2115&-5&2117&8\\
2119&8&2121&-3&2122&-3&2124&8\\
2126&4&2129&5&2130&-7&2131&8\\
2132&8&2133&8&2137&8&2139&4\\
2140&-7&2145&5&2148&-6&2150&-4\\
2151&8&2155&8&2157&8&2158&-4\\
2162&-3&2163&-4&2164&3&2165&-4\\
2168&8&2172&4&2174&4&2177&-4\\
2181&-3&2188&2&2190&-6&2195&8\\
2196&-1&2198&-4&2200&8&2209&-2\\
2212&5&2214&-6&2217&4&2219&4\\
2221&-4&2226&-2&2227&4&2228&-1\\
2229&-4&2232&-1&2238&-4&2242&-7\\
2244&-2&2260&7&2261&8&2268&4\\
2276&-6&2278&-4&2310&1&2314&-5\\
2332&3&2348&8&2360&-2&2379&8\\
\end{array}
\]
}

{\small
\[
\begin{array}{rr@{\qquad}rr@{\qquad}rr@{\qquad}rr}
r&e_r&r&e_r&r&e_r&r&e_r\\ \hline
2380&6&2382&4&2388&2&2390&-4\\
2404&-2&2406&8&2412&8&2416&1\\
2424&4&2458&8&2470&4&2512&5\\
2544&8&2616&6&2672&7&2800&4\\
2816&6&3184&4&3248&8&3264&2\\
3376&-4&3472&4&3520&-4&3552&4\\
3840&-2&4354&3&4355&-4&4357&8\\
4358&8&4364&-4&4614&-3&4615&-1\\
4617&-6&4619&1&4620&4&4622&4\\
4630&3&5132&4&5134&-6&5135&4\\
5138&8&5139&-7&5140&2&5148&4\\
5150&4&5153&-2&5154&2&5155&7\\
5170&4&5172&-6&5190&-3&5204&-6\\
5212&8&5240&-6&5304&4&6168&-4\\
6174&8&6175&-4&6180&-4&6182&3\\
6183&-4&6185&-6&6186&3&6188&4\\
6192&4&6193&6&6194&-4&6195&4\\
6196&4&6197&4&6201&-4&6204&8\\
6206&-4&6211&-4&6213&1&6214&-3\\
6215&4&6217&-2&6218&-5&6219&4\\
6220&2&6222&-4&6225&3&6226&-3\\
6227&4&6228&2&6233&8&6244&-4\\
6258&8&6260&4&6268&4&6275&8\\
6277&1&6284&2&6286&4&6292&4\\
6294&8&6305&-1&6308&1&6310&-4\\
6322&4&6324&-4&6328&2&6338&-2\\
6348&4&6356&8&6392&-2&6410&7\\
6428&-6&6476&-6&6484&2&6500&-6\\
6912&4&7280&4&7344&4&7472&-4\\
\end{array}
\]
}

\begin{corollary}\label{cor:weight-seven-nonextension}
The character $\eta_7$ has no $G$-invariant continuous extension to
$D_3$, $N_{\mathrm{rad}}$ or $N$.
Moreover, writing $\psi_q=\rho_q^I/2\pmod2$ for the divided
odd-prime Kummer characters on $N$, one still has
\[
 \ker\bigl[H^1(N,\mathbf F_2)^G\longrightarrow
                         H^1(D_7,\mathbf F_2)\bigr]
 =\left(\bigoplus_{q\text{ odd prime}}\mathbf F_2\psi_q\right)
       \oplus\mathbf F_2\eta_3.
\]
On $D_3$, the corresponding kernel is exactly $\mathbf F_2\eta_5$.
\end{corollary}

\begin{proof}
Use the arithmetic elements constructed above, for which
\[
 z_3\in J_3,\qquad
 z_5=h z_3h^{-1}z_3^{-t^5}\in J_5,
 \qquad \eta_7(z_5)=1.
\]
For every continuous character $\chi$ to $\mathbf F_2$ and every
profinite exponent $c$, continuity of powers gives
$\chi(z^c)=(c\bmod2)\chi(z)$. If $\chi$ is $G$-invariant on
$D_3$, then
\[
 \chi(z_5)=\chi(z_3)-(t_2^5\bmod2)\chi(z_3)=0,
\]
because $t_2$ is odd. Thus its restriction to $D_5$ cannot be
$\eta_7$. Restricting an invariant character from either larger
group to $D_3$ proves the other nonextension assertions.

Now suppose that a $G$-invariant character of $N$ vanishes on
$D_7$. Its restriction to $D_5$ factors through
$D_5/D_7=128\mathbf Z_2$, so is zero or $\eta_7$.
Nonextension excludes the latter, and hence it vanishes on $D_5$.
Its restriction to $D_3$ is similarly zero or $\eta_5$.
The weight-five arithmetic word
$z_3=h z_1h^{-1}z_1^{-t^3}$, with $z_1\in N_{\mathrm{rad}}$,
satisfies $\eta_5(z_3)=1$, whereas the same odd-power calculation
makes every invariant character of $N$ zero on $z_3$.
Consequently the character vanishes on $D_3$.
The established cubic joint quotient
\[
 N/D_3\simeq
 \left(\prod_{q\text{ odd prime}}2\widehat{\mathbf Z}\right)
       \times8\mathbf Z_2
\]
now gives exactly the displayed direct sum. Conversely every one
of those characters already vanishes on $D_3$ and hence on $D_7$.
For a character defined on $D_3$, the first nonextension argument
again forces vanishing on $D_5$; the quotient
$D_3/D_5=32\mathbf Z_2$ then gives precisely $\mathbf F_2\eta_5$.
\end{proof}

If $J_5=G_{L_5}$ and $L_7=\overline{\mathbf Q}^{\,J_7}$, the exact
image also gives $\operatorname{Gal}(L_7/L_5)\simeq\mathbf Z_2$.
This specifies the extension by the weight-seven monodromy
coordinate; no expression by rational radicals is asserted.
